\documentclass[11pt,a4paper]{article}
\usepackage[T1]{fontenc}
\usepackage[utf8]{inputenc}
\usepackage{amsmath,amsthm,mathtools,aliascnt}
\usepackage{newtxtext,newtxmath}
\usepackage[margin=27mm,headheight=15pt]{geometry}
\usepackage{microtype}
\usepackage{booktabs,array,longtable}
\usepackage[dvipsnames]{xcolor}
\usepackage{enumitem}
\usepackage{fancyhdr}
\usepackage{titlesec}
\usepackage{listings}
\usepackage[most]{tcolorbox}
\usepackage{hyperref}
\usepackage[nameinlink,noabbrev]{cleveref}
\definecolor{OliveInk}{HTML}{53624B}
\definecolor{SoftOlive}{HTML}{EFF2E9}
\definecolor{WarmGray}{HTML}{69675F}
\definecolor{RuleGray}{HTML}{CED3C4}
\hypersetup{colorlinks=true,linkcolor=OliveInk,citecolor=OliveInk,urlcolor=OliveInk,
  pdftitle={A cubic supercongruence for OEIS A348410},
  pdfauthor={ChatGPT; prepared for Vladimir Reshetnikov},
  pdfsubject={An elementary proof of Bala's recorded supercongruence and a two-parameter extension},
  pdfkeywords={OEIS A348410, supercongruence, prime powers, generating functions, reciprocal squares}}
\urlstyle{same}
\titleformat{\section}{\large\bfseries\color{OliveInk}}{\thesection}{0.7em}{}
\titleformat{\subsection}{\normalsize\bfseries\color{OliveInk}}{\thesubsection}{0.7em}{}
\titlespacing*{\section}{0pt}{2.5ex plus 0.8ex minus 0.2ex}{1.0ex}
\titlespacing*{\subsection}{0pt}{1.8ex plus 0.5ex}{0.7ex}
\pagestyle{fancy}
\fancyhf{}
\fancyhead[L]{\small\color{WarmGray}A cubic supercongruence for A348410}
\fancyhead[R]{\small\color{WarmGray}Research note}
\fancyfoot[C]{\small\color{WarmGray}\thepage}
\renewcommand{\headrulewidth}{0.3pt}
\renewcommand{\footrulewidth}{0pt}
\setlength{\parindent}{1.15em}
\setlength{\parskip}{0.25em}
\setlength{\emergencystretch}{1.5em}
\setcounter{tocdepth}{1}
\numberwithin{equation}{section}
\newtheorem{theorem}{Theorem}[section]
\newaliascnt{lemma}{theorem}
\newtheorem{lemma}[lemma]{Lemma}
\aliascntresetthe{lemma}
\newaliascnt{corollary}{theorem}
\newtheorem{corollary}[corollary]{Corollary}
\aliascntresetthe{corollary}
\newaliascnt{proposition}{theorem}
\newtheorem{proposition}[proposition]{Proposition}
\aliascntresetthe{proposition}
\theoremstyle{definition}
\newaliascnt{definition}{theorem}
\newtheorem{definition}[definition]{Definition}
\aliascntresetthe{definition}
\newaliascnt{example}{theorem}
\newtheorem{example}[example]{Example}
\aliascntresetthe{example}
\theoremstyle{remark}
\newaliascnt{remark}{theorem}
\newtheorem{remark}[remark]{Remark}
\aliascntresetthe{remark}
\crefname{theorem}{Theorem}{Theorems}
\crefname{lemma}{Lemma}{Lemmas}
\crefname{corollary}{Corollary}{Corollaries}
\crefname{proposition}{Proposition}{Propositions}
\crefname{remark}{Remark}{Remarks}
\newcommand{\Z}{\mathbb Z}
\newcommand{\Q}{\mathbb Q}
\newcommand{\Zp}{\Z_{(p)}}
\newcommand{\vp}{v_p}
\newcommand{\epsp}{\varepsilon_p}
\newcommand{\coeff}[2]{[x^{#1}]\,#2}
\newcommand{\Auv}{A_{u,v}}
\newcommand{\Fp}{F_{u,v}}
\newcommand{\Sp}{S_p}
\newcommand{\Hp}{H_p}
\newcommand{\Cp}{C_p}
\newcommand{\dd}{\,\mathrm d}
\newcommand{\pmodring}[1]{\pmod{#1\Zp}}
\lstset{basicstyle=\small\ttfamily,keywordstyle=\color{OliveInk}\bfseries,
  commentstyle=\color{WarmGray},stringstyle=\color{OliveInk},
  showstringspaces=false,breaklines=true,columns=fullflexible,
  frame=single,rulecolor=\color{RuleGray},backgroundcolor=\color{SoftOlive},
  xleftmargin=0.5em,xrightmargin=0.5em,framesep=0.65em}

\begin{document}
\thispagestyle{empty}
\begin{center}
{\small\scshape\color{OliveInk} Integer sequences \quad / \quad Prime-power congruences}\par
\vspace{1.2em}
{\LARGE\bfseries A cubic supercongruence\par for OEIS A348410}\par
\vspace{0.8em}
{\large An elementary proof, a two-parameter extension,\par and a boundary counterexample}\par
\vspace{1.0em}
{\normalsize Research note prepared for Vladimir Reshetnikov}\par
{\small 19 September 2026}\par
\end{center}
\vspace{0.7em}
\begin{abstract}
Let $a(0)=1$ and
\[
 a(n)=[x^n]\bigl((1-x)(1-x^2)\bigr)^{-n}.
\]
We give a self-contained proof of the conjecture recorded by Peter Bala in
OEIS A348410: for every prime $p\geq5$ and all positive integers $m,r$,
\[
 a(mp^r)\equiv a(mp^{r-1})\pmod{p^{3r}}.
\]
In fact, the same congruence holds for
$A_{u,v}(n)=[x^n](1-x)^{-un}(1+x)^{-vn}$ with arbitrary integers $u,v$.
At $p=3$, the uniform exponent becomes $3r-1$; its loss of one power is
necessary already for the original sequence. The proof splits a formal
logarithm into degrees divisible and not divisible by $p$. After an
exponential truncation, only a quadratic coefficient can matter. Its
extra divisibility follows from a reciprocal-square sum and a parity
identity; a logarithmic-derivative estimate then completes the argument
at every prime-power level. We also derive a denominator restriction for
the cubic M\"obius transform and exhibit a period-three replacement that
fails modulo $5^3$. A separate source audit explains why a broader printed
framing claim is not used. Exact-arithmetic programs, test data, and
symbolic generating-function certificates accompany the note.
\end{abstract}
\begin{tcolorbox}[colback=SoftOlive,colframe=RuleGray,boxrule=0.4pt,arc=1mm,
  left=2mm,right=2mm,top=1.2mm,bottom=1.2mm]
\small\textbf{Scope of the result.}
The argument below proves the precise conjecture still displayed in the
OEIS entry inspected on the date above. The generating function and
recurrence literature is credited separately. This is a research
manuscript, not a claim of independent refereeing or established
historical priority.
\end{tcolorbox}
\vfill
{\small\color{WarmGray}Companion files: \texttt{code/verify.py}, exact CSV data,
optional symbolic checks, and a proposed OEIS annotation.}
\newpage
\tableofcontents
\vspace{1.0em}

\section{The problem and the theorem}\label{sec:problem}

For $n\geq1$, consider $n$ labelled pairs of baskets. The first basket in
each pair may contain any nonnegative number of indistinguishable
objects, and the second must contain an even number. The generating
function for one pair is
\begin{equation}\label{eq:F-original}
 F(x)=\frac{1}{(1-x)(1-x^2)}
     =\frac{1}{(1-x)^2(1+x)}.
\end{equation}
The number of ways to distribute exactly $n$ objects among $n$ such pairs
is therefore $a(n)=[x^n]F(x)^n$. This is the definition of OEIS A348410
\cite{oeis}. Its first terms, including $a(0)=1$, are
\[
1,1,5,19,85,376,1715,7890,36693,171820,809380,\ldots.
\]
The entry attributes to Peter Bala, on 21 February 2022, the assertion
\begin{equation}\label{eq:bala}
 a(mp^r)\equiv a(mp^{r-1})\pmod{p^{3r}}
 \qquad(p\geq5\text{ prime};\ m,r\geq1)
\end{equation}
as a conjecture \cite{oeis}. The goal here is to prove that assertion,
not merely to extend its numerical verification.

\subsection{Why this particular target}

The displayed coefficient formula is much more useful for divisibility
than a high-order recurrence in $n$. Its logarithmic derivative has a
coefficient sequence of period two, and multiplication of an index by an
odd prime preserves its parity. That observation gives an exact
factorization rather than an approximate congruence. The nontrivial
remaining issue is whether a weighted reciprocal-square sum supplies
one more full prime-power factor. We prove that it does in this setting.

Recent papers by Niu \cite{niu} and Prodinger \cite{prodinger} study the
algebraic ordinary generating function; Prodinger also extracts a
recurrence. Those are not claimed as new here. A derivation of the known
parametric generating function is included in \cref{sec:gf} to connect
the arithmetic proof with the combinatorial sequence. Searches centered
on the sequence identifier and its supercongruence did not locate a
direct proof of \eqref{eq:bala}. This is a bounded literature observation,
not a proof that no prior treatment exists under other notation.

\subsection{The larger family}

For $u,v\in\Z$, set
\begin{equation}\label{eq:family}
 \Fp(x)=(1-x)^{-u}(1+x)^{-v},\qquad
 \Auv(n)=[x^n]\Fp(x)^n,\quad n\geq0.
\end{equation}
Integer powers, including negative ones, are understood as formal power
series at $x=0$. Their coefficients are integral; hence
$\Auv(n)\in\Z$. In particular, $\Auv(0)=1$ and $a(n)=A_{2,1}(n)$.
For an odd prime define
\begin{equation}\label{eq:epsilon}
 \epsp=\begin{cases}1,&p=3,\\0,&p\geq5.\end{cases}
\end{equation}

\begin{theorem}[Prime-power congruence]\label{thm:main}
Let $u,v\in\Z$, let $p$ be an odd prime, and let $N$ be a positive integer
with $p\mid N$. Write $r=\vp(N)$. Then
\begin{equation}\label{eq:main}
 \Auv(N)-\Auv(N/p)\in p^{3r-\epsp}\Z.
\end{equation}
Equivalently, for all $m,k\geq1$,
\begin{equation}\label{eq:strong-m}
 \Auv(mp^k)\equiv\Auv(mp^{k-1})
 \pmod{p^{3(k+\vp(m))-\epsp}}.
\end{equation}
\end{theorem}

Taking $(u,v)=(2,1)$ and $p\geq5$ proves \eqref{eq:bala}, including the
case $p\mid m$. The formulation with the actual valuation of $N$ avoids
a common unnecessary restriction that the multiplier be prime to $p$.
The proof occupies \cref{sec:tools,sec:harmonic,sec:proof}.

\section{Formal series and coefficient divisibility}\label{sec:tools}

\subsection{The arithmetic setting}

Fix an odd prime $p$. We work in the localization
\[
 \Zp=\{a/b\in\Q:p\nmid b\}.
\]
Write $\vp(0)=+\infty$, and extend the usual valuation to nonzero rational
numbers by $\vp(a/b)=\vp(a)-\vp(b)$. A congruence between formal power
series over $\Zp$ is always coefficientwise. For example,
$U\equiv V\pmod{p^s}$ means every coefficient of $U-V$ lies in
$p^s\Zp$.

All logarithms and exponentials below are formal series over $\Q$.
Their use requires no analytic convergence: when a series $H$ has zero
constant term, the coefficient of $x^d$ in $\exp(H)$ receives
contributions from only $H^k$ with $k\leq d$. Valuation estimates are
then applied to these finite coefficient sums.

\begin{lemma}[Logarithmic-derivative estimate]\label{lem:derivative}
Let $G(x)\in1+x\Zp[[x]]$ and write
\[
 x\frac{G'(x)}{G(x)}=\sum_{i\geq1}b_i x^i,
 \qquad b_i\in\Zp.
\]
For a positive integer $L$, put $y_j=[x^j]G(x)^L$. Then, for $j\geq1$,
\begin{equation}\label{eq:derivative-bound}
 \vp(y_j)\geq\max\{0,\vp(L)-\vp(j)\}.
\end{equation}
Also $y_0=1$.
\end{lemma}
\begin{proof}
Taking coefficients in
$x(G^L)'=L(xG'/G)G^L$ gives the exact identity
\begin{equation}\label{eq:derivative-coeff}
 j y_j=L\sum_{i=1}^{j}b_i y_{j-i}.
\end{equation}
The sum on the right belongs to $\Zp$. Thus
$\vp(y_j)\geq\vp(L)-\vp(j)$. Independently, $G^L$ has coefficients in
$\Zp$, so $\vp(y_j)\geq0$. Combining the two bounds proves the result.
\end{proof}

The estimate is particularly valuable when a complementary index is
more highly divisible by $p$ than the index $j$ itself. No asymptotics or
induction on a guessed recurrence is involved.

\begin{lemma}[Exponential truncation]\label{lem:tail}
Suppose $r=\vp(N)\geq1$ and $H(x)\in x\Zp[[x]]$. Then
\begin{equation}\label{eq:tail}
 \exp(NH)\equiv1+NH+\frac{N^2}{2}H^2
 \pmod{p^{3r-\epsp}\Zp[[x]]}.
\end{equation}
\end{lemma}
\begin{proof}
For every $k\geq3$, each coefficient of $(N^k/k!)H^k$ has valuation at
least $kr-\vp(k!)$. The elementary factorial estimate
\[
 \vp(k!)=\sum_{j\geq1}\left\lfloor\frac{k}{p^j}\right\rfloor
 \leq\frac{k-1}{p-1}
\]
follows, for example, by writing $k$ in base $p$.
For $p\geq5$, it implies $\vp(k!)\leq k-3$ when $k\geq4$;
the same bound holds directly at $k=3$. Hence
\[
 kr-\vp(k!)\geq3r+(k-3)(r-1)\geq3r.
\]
For $p=3$, the bound $\vp(k!)\leq(k-1)/2\leq k-2$ gives
\[
 kr-\vp(k!)\geq3r-1+(k-3)(r-1)\geq3r-1.
\]
All omitted terms consequently lie in the stated ideal, coefficient by
coefficient.
\end{proof}

\subsection{An exact separation of the logarithm}

For $G=\Fp$, direct expansion gives
\begin{equation}\label{eq:log}
 H(x):=\log\Fp(x)=\sum_{j\geq1}\frac{b_j}{j}x^j,
 \qquad b_j=u+v(-1)^j.
\end{equation}
In the original case, $b_j=1$ for odd $j$ and $b_j=3$ for even $j$.
For every odd prime, $b_{pj}=b_j$. Define the part with degrees not
divisible by $p$ by
\begin{equation}\label{eq:Hp}
 \Hp(x)=\sum_{\substack{j\geq1\\p\nmid j}}\frac{b_j}{j}x^j
 \in x\Zp[[x]].
\end{equation}
Separating the multiples of $p$ in \eqref{eq:log} yields
\begin{equation}\label{eq:logsplit}
 H(x)=\frac1pH(x^p)+\Hp(x).
\end{equation}
Therefore, whenever $L=N/p$ is an integer,
\begin{equation}\label{eq:exact-factorization}
 \Fp(x)^N=\Fp(x^p)^L\exp\bigl(N\Hp(x)\bigr).
\end{equation}
This is an identity over $\Q[[x]]$, not an application of a congruence
whose error must be estimated.

The linear term $N\Hp$ in \eqref{eq:tail} cannot contribute to the
coefficient of $x^N$ in \eqref{eq:exact-factorization}: $N$ is divisible
by $p$, the first factor has only degrees divisible by $p$, and $\Hp$
has none. The constant term gives $\Auv(L)$. Thus the entire problem
reduces to controlling the quadratic term.

\section{The reciprocal-square cancellation}\label{sec:harmonic}

\subsection{The unweighted sum}

For a positive integer $T$ divisible by $p$, define
\begin{equation}\label{eq:S}
 \Sp(T)=\sum_{\substack{1\leq i<T\\p\nmid i}}\frac1{i^2}\in\Zp.
\end{equation}

\begin{lemma}[Reciprocal squares over complete blocks]\label{lem:squares}
For every odd prime $p$ and $T$ with $p\mid T$,
\begin{equation}\label{eq:S-bound}
 \vp\bigl(\Sp(T)\bigr)\geq\vp(T)-\epsp.
\end{equation}
\end{lemma}
\begin{proof}
Write $s=\vp(T)$ and $T=q p^s$. Reduction modulo $p^s$ partitions the
sum into $q$ copies of the nonzero unit residues modulo $p^s$. Thus
\[
 \Sp(T)\equiv q U_s\pmodring{p^s},\qquad
 U_s=\sum_{\substack{1\leq a<p^s\\p\nmid a}}a^{-2}.
\]
Multiplication by $2$ permutes these unit residue classes. Consequently,
\[
 U_s\equiv\sum_{a\in(\Z/p^s\Z)^\times}(2a)^{-2}
     \equiv\frac14U_s\pmodring{p^s}.
\]
It follows that $3U_s\equiv0\pmod{p^s}$. Since
$\vp(3)=\epsp$, we obtain $\vp(U_s)\geq s-\epsp$ and hence the same
bound for $\Sp(T)$.
\end{proof}

The factor $3$ makes the exception at $p=3$ transparent. The proof also
shows why an unweighted sum is safe to permute: the summand depends
only on the unit residue modulo $p^s$.

\subsection{Parity weights}

When $M$ is even and divisible by the odd prime $p$, separating its even
indices gives the exact identity
\begin{equation}\label{eq:alternating}
 \sum_{\substack{1\leq i<M\\p\nmid i}}\frac{(-1)^i}{i^2}
 =\frac12\Sp(M/2)-\Sp(M).
\end{equation}
Indeed, the sum over even indices is $\tfrac14\Sp(M/2)$, and the
alternating sum is twice that sum minus the full sum. Because
$\vp(M/2)=\vp(M)$, \cref{lem:squares} gives the same lower bound for
both terms on the right.

\begin{lemma}[Quadratic coefficient estimate]\label{lem:quadratic}
Let $\Hp$ be defined by \eqref{eq:Hp}, with $b_j=u+v(-1)^j$.
For every positive multiple $M$ of $p$, set
\[
 \Cp(M)=[x^M]\Hp(x)^2.
\]
Then
\begin{equation}\label{eq:C-bound}
 \vp\bigl(\Cp(M)\bigr)\geq\vp(M)-\epsp.
\end{equation}
\end{lemma}
\begin{proof}
Put $s=\vp(M)$. Since $p\mid M$, the conditions $p\nmid i$ and
$p\nmid(M-i)$ are equivalent, and
\begin{equation}\label{eq:C-expanded}
 \Cp(M)=\sum_{\substack{1\leq i<M\\p\nmid i}}
 \frac{b_i b_{M-i}}{i(M-i)}.
\end{equation}
The exact identity
\[
 \frac1{i(M-i)}+\frac1{i^2}=\frac{M}{i^2(M-i)}
\]
shows that
\begin{equation}\label{eq:C-reduced}
 \Cp(M)\equiv-
 \sum_{\substack{1\leq i<M\\p\nmid i}}\frac{b_i b_{M-i}}{i^2}
 \pmodring{p^s}.
\end{equation}
There are two parity cases.

If $M$ is odd, $i$ and $M-i$ have opposite parity, so
$b_i b_{M-i}=u^2-v^2$ is independent of $i$. Thus
\begin{equation}\label{eq:C-odd}
 \Cp(M)\equiv-(u^2-v^2)\Sp(M)\pmodring{p^s}.
\end{equation}
The desired bound follows from \cref{lem:squares}.

If $M$ is even, the two indices have the same parity and
\[
 b_i b_{M-i}=u^2+v^2+2uv(-1)^i.
\]
Substituting \eqref{eq:alternating} into \eqref{eq:C-reduced} simplifies
the weighted sum to
\begin{equation}\label{eq:C-even}
 \Cp(M)\equiv-(u-v)^2\Sp(M)-uv\Sp(M/2)
 \pmodring{p^s}.
\end{equation}
Both terms on the right have valuation at least $s-\epsp$.
The error term in either parity case has valuation at least $s$, which
is no smaller. This proves \eqref{eq:C-bound}.
\end{proof}

The parity reduction is the essential structural step. Periodicity
alone would not justify replacing a weighted sum by a permuted one:
reducing an index modulo $p^s$ can change its residue modulo the period.
The exact decomposition \eqref{eq:alternating}, rather than such a
permutation of the weights, is what makes the argument valid.

\section{Proof of the prime-power theorem}\label{sec:proof}

\begin{proof}[Proof of \cref{thm:main}]
Fix $N$, put $r=\vp(N)\geq1$, and write $L=N/p$. To shorten notation,
write $G=\Fp$ and
\[
 y_j=[x^j]G(x)^L.
\]
By the exact factorization \eqref{eq:exact-factorization} and the
truncation \eqref{eq:tail},
\begin{equation}\label{eq:quadratic-reduction}
 \Auv(N)-\Auv(L)
 \equiv\frac{N^2}{2}[x^N]\,G(x^p)^L\Hp(x)^2
 \pmodring{p^{3r-\epsp}}.
\end{equation}
As explained above, the constant term supplies $\Auv(L)$ and the linear
term supplies zero. Expanding the coefficient on the right gives
\begin{equation}\label{eq:convolution}
 [x^N]G(x^p)^L\Hp(x)^2
   =\sum_{m=1}^{L}y_{L-m}\Cp(pm).
\end{equation}
The index $m=0$ is absent because $\Hp^2$ has zero constant term.
We claim that every summand in \eqref{eq:convolution} belongs to
$p^{r-\epsp}\Zp$.

Suppose first that $t=\vp(m)<r-1$. Since $\vp(L)=r-1$, the two terms
$L$ and $m$ have different valuations, so
\[
 \vp(L-m)=t.
\]
In particular $L-m>0$. By \cref{lem:derivative,lem:quadratic},
\begin{align}
 \vp(y_{L-m})&\geq r-1-t,\label{eq:first-factor}\\
 \vp\bigl(\Cp(pm)\bigr)&\geq1+t-\epsp.\label{eq:second-factor}
\end{align}
Their sum is $r-\epsp$.

Suppose instead that $\vp(m)\geq r-1$. In this case,
\[
 \vp\bigl(\Cp(pm)\bigr)\geq1+\vp(m)-\epsp\geq r-\epsp,
\]
while $y_{L-m}$ is integral at $p$. This also handles the endpoint
$m=L$, where $y_0=1$ and no valuation of the zero index is needed.
The claim is proved in all cases, including $r=1$, for which the first
case is empty.

It follows from \eqref{eq:convolution} that the quadratic coefficient has
valuation at least $r-\epsp$. Since $p$ is odd,
$\vp(N^2/2)=2r$. The right side of \eqref{eq:quadratic-reduction}
therefore belongs to $p^{3r-\epsp}\Zp$. Its left side is an integer, so
it lies in $p^{3r-\epsp}\Z$, as required. Finally, substituting
$N=mp^k$ gives \eqref{eq:strong-m}.
\end{proof}

\subsection{Where the three powers come from}

The exponent is not a numerical coincidence. Two copies of $N$ come
from the quadratic exponential term. For a term indexed by $m$, the
reciprocal-square estimate contributes $1+\vp(m)$ powers of $p$, except
for the possible loss at $3$. When that contribution is less than $r$,
the coefficient $y_{L-m}$ contributes exactly the amount needed to
reach $r$. These complementary estimates explain why a proof only for
$N=p$ would miss an important part of the argument.

\begin{remark}[A sufficient criterion beyond this family]\label{rem:criterion}
At a fixed prime $p\geq5$, the same proof applies to any
$G\in1+x\Zp[[x]]$ for which the coefficients
$xG'/G=\sum_{j\geq1}b_jx^j$ satisfy $b_{pj}=b_j$, provided that
\[
 [x^M]\left(\sum_{p\nmid j}\frac{b_j}{j}x^j\right)^2
 \in p^{\vp(M)}\Zp\qquad(p\mid M).
\]
Thus exact invariance under index multiplication and a quadratic
coefficient estimate are sufficient. The latter is an additional
hypothesis, not a consequence of arbitrary periodicity.
\end{remark}

\section{Consequences and sharp examples}\label{sec:consequences}

\subsection{The original conjecture and the small primes}

Theorem~\ref{thm:main} immediately gives the full conjecture for
$A_{2,1}=a$. It also proves
\begin{equation}\label{eq:three-adic}
 a(m3^r)\equiv a(m3^{r-1})\pmod{3^{3r-1}},\qquad m,r\geq1.
\end{equation}
Neither the uniform improvement $3r\mapsto3r+1$ for all $p\geq5$ nor
the replacement $3r-1\mapsto3r$ at $p=3$ is possible. Indeed,
\begin{align*}
 a(5)-a(1)&=376-1=375=3\cdot5^3,\\
 a(3)-a(1)&=19-1=18=2\cdot3^2.
\end{align*}
These are sharpness statements for the uniform formulas; they do not
assert equality of the valuation for every prime, multiplier, or level.
At $p=2$, even the cubic congruence at the first level fails:
$a(2)-a(1)=4$ is not divisible by $2^3$. The odd-prime theorem makes no
claim about an optimal $2$-adic replacement.

For illustration, exact values along several prime-power towers are
shown in \cref{tab:valuations}. Every entry is computed from integers,
not inferred from a decimal approximation.
\begin{table}[htbp]
\centering
\caption{Observed valuations of $a(p^r)-a(p^{r-1})$; $m=1$.}
\label{tab:valuations}
\begin{tabular}{c r r r r}
\toprule
$p$ & $r=1$ & $r=2$ & $r=3$ & $r=4$\\
\midrule
3 & 2 & 5 & 8 & 11\\
5 & 3 & 6 & 9 & 12\\
7 & 3 & 6 & 9 & 12\\
11 & 3 & 6 & 9 & --\\
\bottomrule
\end{tabular}
\end{table}
The dash means the index $11^4$ lies outside the default computational
sweep, not that the congruence is unknown there. As a larger explicit
check,
\[
 a(25)-a(5)=13\,859\,743\,127\,937\,500
\]
has $5$-adic valuation $6$.

\subsection{Other integer-sequence families}

Writing $u=\alpha+\beta$ and $v=\beta$ gives
\begin{equation}\label{eq:alpha-beta}
 A_{\alpha+\beta,\beta}(n)
   =[x^n](1-x)^{-\alpha n}(1-x^2)^{-\beta n}.
\end{equation}
For nonnegative integers $\alpha,\beta$, this counts distributions into
$\alpha n$ unrestricted baskets and $\beta n$ even-occupancy baskets.
All these sequences satisfy the same odd-prime congruences. Negative
values of $u,v$ are also covered algebraically, although the positive
basket interpretation then need not apply.

Several elementary specializations help check the formulation.
For a positive integer $u$,
\[
 A_{u,0}(n)=\binom{(u+1)n-1}{n}\qquad(n\geq1).
\]
For a positive integer $q$, $F_{q,q}(x)=(1-x^2)^{-q}$, so
\[
 A_{q,q}(n)=
 \begin{cases}
 0,&n\text{ odd},\\
 \displaystyle\binom{qn+n/2-1}{n/2},&n\text{ even}.
 \end{cases}
\]
The theorem accommodates both the zero odd terms and the nonzero even
terms without a separate divisibility convention.

There is also a simple change-of-step corollary. Let $d\geq1$ and put
\[
 B(n)=[x^n](1-x^d)^{-un}(1+x^d)^{-vn}.
\]
For a prime $p\nmid d$, the conditions $d\mid N$ and $d\mid N/p$ are
equivalent. If they fail, both relevant coefficients vanish. If they
hold, $B(N)=A_{du,dv}(N/d)$, and
$\vp(N/d)=\vp(N)$. Thus the same bound holds for $B$ at every odd prime
not dividing $d$.

\subsection{A cubic M\"obius transform}

Let $\mu$ denote the M\"obius function. Define the rational numbers
\begin{equation}\label{eq:mobius}
 c_{u,v}(n)=\frac1{n^3}\sum_{d\mid n}\mu(n/d)\Auv(d),\qquad n\geq1.
\end{equation}
These are arithmetic transforms, not automatically counts of primitive
combinatorial objects.

\begin{corollary}[Denominators of the cubic transform]\label{cor:mobius}
For every $u,v\in\Z$ and $n\geq1$,
\begin{equation}\label{eq:denominators}
 c_{u,v}(n)\in\frac13\Z[1/2].
\end{equation}
In particular, $3c_{u,v}(n)$ is an integer when $n$ is odd, and
$c_{u,v}(n)$ is an integer when $\gcd(n,6)=1$.
\end{corollary}
\begin{proof}
Fix an odd prime $p\mid n$ and write $n=p^s w$ with $p\nmid w$.
Only divisors with $p$-exponent $s$ or $s-1$ survive the M\"obius factor.
Consequently the numerator in \eqref{eq:mobius} equals
\begin{equation}\label{eq:mobius-pairing}
 \sum_{e\mid w}\mu(w/e)
 \bigl(\Auv(p^s e)-\Auv(p^{s-1}e)\bigr).
\end{equation}
Each difference is divisible by $p^{3s-\epsp}$ by
\cref{thm:main}. Dividing by $n^3$ shows
$\vp(c_{u,v}(n))\geq-\epsp$. No denominator prime can occur unless it
divides $n$, because the numerator in \eqref{eq:mobius} is an integer.
Hence primes $p\geq5$ never occur in the reduced denominator, and the
power of $3$ is at most one. This is exactly \eqref{eq:denominators};
the two stated special cases follow.
\end{proof}

For the original sequence, the values of $c_{2,1}(n)$ for $1\leq n\leq12$
are
\[
 1,\ \frac12,\ \frac23,\ \frac54,\ 3,\ \frac{47}{6},\ 23,
 \ \frac{143}{2},\ \frac{707}{3},\ 809,\ 2878,\ \frac{126385}{12}.
\]
M\"obius inversion also supplies the formal Lambert-series identity
\begin{equation}\label{eq:lambert}
 \sum_{n\geq1}\Auv(n)t^n
  =\sum_{d\geq1}d^3c_{u,v}(d)\frac{t^d}{1-t^d}.
\end{equation}
Thus the supercongruence has an equivalent denominator interpretation
in a generating-function transform.

\subsection{Limits along a prime-power tower}

For a fixed odd prime $p$ and positive integer $m$, the sequence
$\Auv(mp^r)$ is Cauchy in $\Z_p$ and therefore has a $p$-adic limit
$\lambda_{u,v,p,m}$. Summing the successive differences from
\eqref{eq:strong-m} gives the error estimate
\[
 \vp\bigl(\lambda_{u,v,p,m}-\Auv(mp^r)\bigr)
 \geq3(r+1+\vp(m))-\epsp\qquad(r\geq0).
\]
This conclusion does not require identifying the limit in closed form.
It expresses another consequence of the full tower result rather than
just the first congruence modulo $p^3$.

\section{Generating functions and exact evaluation}\label{sec:gf}

This section places the arithmetic theorem alongside the ordinary
generating-function description. The parametric and quartic formulas
are known from the sequence entry and the recent treatments
\cite{oeis,niu,prodinger}; the derivations here are included for context
and to make the computational checks transparent.

\subsection{Coefficient sums}

Expanding the two factors in \eqref{eq:F-original} separately gives, for
$n\geq1$,
\begin{equation}\label{eq:positive-sum}
 a(n)=\sum_{k=0}^{\lfloor n/2\rfloor}
 \binom{2n-2k-1}{n-2k}\binom{n+k-1}{k}.
\end{equation}
Every term is nonnegative. For example, at $n=5$ the three summands are
$126$, $175$, and $75$, giving $a(5)=376$.
This formula and the coefficient recurrence below provide independent
ways of obtaining the same integers.

\subsection{A parametric ordinary generating function}

Let $y(t)$ be the unique formal series with $y(0)=0$ satisfying
\begin{equation}\label{eq:y}
 y=tF(y),\qquad\text{or equivalently}\qquad
 t=y(1-y)^2(1+y)=y-y^2-y^3+y^4.
\end{equation}
The linear coefficient in the last expression is $1$, so this formal
inverse exists and is unique. Lagrange inversion, applied to $\log F$,
gives for $n\geq1$
\begin{align*}
 [t^n]\log\frac{y(t)}t
 &=\frac1n[x^{n-1}]\frac{F'(x)}{F(x)}F(x)^n\\
 &=\frac1{n^2}[x^{n-1}]\bigl(F(x)^n\bigr)'\\
 &=\frac1n[x^n]F(x)^n=\frac{a(n)}n.
\end{align*}
Consequently, with $\mathcal A(t)=\sum_{n\geq0}a(n)t^n$,
\[
 \mathcal A(t)=1+t\frac{\mathrm d}{\mathrm dt}\log\frac{y(t)}t
             =\frac{t y'(t)}{y(t)}.
\]
Differentiating $y=tF(y)$ now yields
\begin{equation}\label{eq:parametric}
 \mathcal A(t)=\frac1{1-yF'(y)/F(y)}
             =\frac{1-y^2}{1-y-4y^2}.
\end{equation}
Eliminating $y$ gives
\begin{align}\label{eq:quartic}
 0={}&(256t^2+107t-32)(\mathcal A^4-\mathcal A^3)\notag\\
    &+(96t^2+36t)\mathcal A^2-(16t^2+4t)\mathcal A+t^2.
\end{align}
The branch is fixed by $\mathcal A(0)=1$. The optional symbolic program
verifies that the resultant of the two defining equations is exactly
the polynomial in \eqref{eq:quartic}. It also checks the identity after
substitution and clearing the fourth power of $1-y-4y^2$.

Algebraicity is not used to prove the supercongruence. The logarithmic
coefficient pattern of $F$, rather than merely the algebraicity of
$\mathcal A$, is the input that supplies the required divisibility.

\subsection{A short exact recurrence for coefficient extraction}

For fixed $n$, write
\[
 f_j=[x^j](1-x)^{-un}(1+x)^{-vn}.
\]
The differential equation
\[
 (1-x^2)\frac{\mathrm d}{\mathrm dx}\Fp(x)^n
   =n\bigl((u-v)+(u+v)x\bigr)\Fp(x)^n
\]
gives
\begin{equation}\label{eq:coefficient-recurrence}
 j f_j=n(u-v)f_{j-1}+\bigl(n(u+v)+j-2\bigr)f_{j-2},
 \quad j\geq1,
\end{equation}
with $f_{-1}=0$, $f_0=1$. The desired value is $f_n$.
This is a recurrence in the coefficient index $j$ with $n$ held fixed;
it should not be confused with a recurrence for the diagonal sequence
as $n$ varies.

Only two previous integers need be retained, and each division by $j$
is exact. There are $O(n)$ arithmetic steps for a single diagonal value,
although the integer bit lengths grow with $n$ and the bit complexity
is not claimed to be linear. In particular, one must not replace the
exact division by a modular inverse when $p\mid j$.

A minimal implementation is:
\begin{lstlisting}[language=Python]
def diagonal(u: int, v: int, n: int) -> int:
    older, previous = 0, 1
    for j in range(1, n + 1):
        numerator = (n * (u - v) * previous
                     + (n * (u + v) + j - 2) * older)
        current, remainder = divmod(numerator, j)
        if remainder:
            raise ArithmeticError("Nonexact coefficient division")
        older, previous = previous, current
    return previous
\end{lstlisting}
The full checker includes input validation, a separate binomial
implementation for signed integer parameters, and an independent
basket-convolution implementation for small indices.

\section{A period-three replacement that fails}\label{sec:boundary}

It is tempting to replace the factor $1-x^2$ by $1-x^d$ and assume that
any fixed period $d$, away from its prime divisors, gives the same cubic
congruence. This is false already for $d=3$.

\begin{proposition}[A small counterexample]\label{prop:counterexample}
Let
\[
 d(n)=[x^n](1-x)^{-n}(1-x^3)^{-n}.
\]
Then $d(1)=1$ and $d(5)=201$. In particular,
\[
 d(5)-d(1)=200=8\cdot5^2\not\equiv0\pmod{5^3}.
\]
\end{proposition}
\begin{proof}
For degree $5$, the factor $(1-x^3)^{-5}$ can contribute only its
constant term or its term $5x^3$. Hence
\[
 d(5)=\binom95+5\binom62=126+75=201.
\]
The remaining assertions follow immediately.
\end{proof}

The failure is not caused by $p$ dividing the period: $5\nmid3$.
The logarithmic derivative has coefficients
\[
 b_j=1+3\,\mathbf1_{3\mid j},
\]
which satisfy $b_{5j}=b_j$. Thus the exact logarithmic splitting still
works at $p=5$. What fails is the extra divisibility of the quadratic
coefficient. Directly,
\begin{align*}
 [x^5]\left(\sum_{5\nmid j}\frac{b_j}{j}x^j\right)^2
 &=2\left(1\cdot\frac14+\frac12\cdot\frac43\right)\\
 &=\frac{11}{6}\not\equiv0\pmod5.
\end{align*}
The quadratic term then contributes two powers of $5$, not necessarily
three. This is precisely the hypothesis isolated in
\cref{rem:criterion}.

There is an infinite parameter family of such first-level failures.
For a positive integer $t$ with $5\nmid t$, define
\[
 d_t(n)=[x^n](1-x)^{-tn}(1-x^3)^{-tn}.
\]
Since
\[
 d_t(5)=\binom{5t+4}{5}+5t\binom{5t+1}{2},\qquad d_t(1)=t,
\]
and \cref{thm:main} with $(u,v)=(t,0)$ gives
$\binom{5t+4}{5}\equiv t\pmod{125}$, one obtains
\begin{equation}\label{eq:period3-family}
 d_t(5)-d_t(1)\equiv\frac{25t^2}{2}\pmod{125}.
\end{equation}
The denominator $2$ is interpreted as a unit modulo $125$. It follows
that $v_5(d_t(5)-d_t(1))=2$ whenever $5\nmid t$.
A member with $t=3$ is useful for the literature audit in
\cref{app:audit}, because its logarithmic coefficient sequence satisfies
the rational $2$-sequence conditions at \emph{every} prime, including $3$.

\section{Verification, limitations, and outcome}\label{sec:verification}

\subsection{What was actually executed}

The default standard-library checker was run with Python 3.13.5 and
\lstinline|--max-n 5000|. It produced the successful checks summarized in
\cref{tab:checks}. The code is written for Python 3.9 or later and does
not make network requests. All arithmetic relevant to a congruence is
integer, rational, or modular arithmetic; no floating-point comparison
is used. The optional polynomial certificate was also run, with
SymPy 1.14.0, and returned zero for the substituted quartic and $1$ for
the quotient of its resultant by the displayed quartic.

\begin{table}[htbp]
\centering
\caption{Finite checks executed by the supplied programs.}
\label{tab:checks}
\small
\begin{tabular}{p{0.68\linewidth}r}
\toprule
Check & Number\\
\midrule
Agreement with the initial values displayed in OEIS & 26\\
Agreement with the positive binomial sum, $0\leq n\leq150$ & 151\\
Independent basket convolution, $0\leq n\leq12$ & 13\\
Signed binomial comparisons for $-3\leq u,v\leq4$, $0\leq n\leq25$ & 1\,664\\
Original-sequence prime-power congruences & 229\\
Two-parameter prime-power congruences & 2\,048\\
Quadratic coefficient bounds & 1\,260\\
Logarithmic-derivative coefficient bounds & 3\,415\\
Factorial inequalities for the exponential tail & 119\,760\\
Coefficients checked in each generating-function identity & 61\\
Rational $2$-sequence input checks for the audit example & 10\,000\\
Cubic M\"obius-transform denominator checks & 100\\
\bottomrule
\end{tabular}
\end{table}

The original-sequence sweep uses all odd primes at most $97$, multipliers
$1\leq m\leq5$, and all positive levels for which $mp^r\leq5000$.
The largest tested index is $4913$. The family sweep uses all $64$ pairs
$-3\leq u,v\leq4$, multipliers $m\in\{1,2,3,5\}$, and levels up to $3$
for $p=3$, up to $2$ for $p=5,7$, and level $1$ for $p=11$.
It checks the stronger bound involving the actual valuation of $mp^r$.
Zero differences are recorded as having infinite valuation.

The derivative and quadratic checks also include signed and degenerate
parameter pairs. Some prescribed bounds are zero and some coefficients
vanish; the counts include those cases rather than silently discarding
them. Complete ranges and all congruence records are recoverable from
the script and CSV files.

\subsection{Reproducing the calculation}

From the archive directory, run
\begin{lstlisting}[language=bash]
python code/verify.py
python code/symbolic_checks.py   # optional; requires SymPy
\end{lstlisting}
The first command rewrites the CSV and JSON files in \texttt{data/}.
The second performs exact polynomial-identity checks; it is not needed
for the supercongruence proof. The TeX source builds with
\texttt{latexmk -pdf article.tex}. There are no external figure files or
font files to install from the archive.

The tests are useful for detecting wrong signs, incorrect exponents,
and off-by-one ranges. They do not establish the universally quantified
statements. Those are proved in the lemmas and the main theorem. No
Lean or other proof-assistant formalization, and no independent referee
verification, is claimed.

\subsection{What is settled and what remains a separate question}

The stated goal is achieved: the congruence recorded in A348410 follows
for every allowed prime, multiplier, and level. The extension to integer
parameters and the $3$-adic version use the same proof. The exact
counterexample rules out the unrestricted replacement of parity by a
larger period.

Several refinements remain distinct from the solved assertion. The
optimal $2$-adic bound for the original sequence is not determined here.
The theorem also does not classify the parameters and primes at which
its lower bound is strict. A systematic classification of logarithmic
coefficient patterns satisfying the quadratic criterion could produce
further families, but any such classification must accommodate the
period-three obstruction rather than assume it away.

\appendix
\section{Audit of a broader framing claim}\label{app:audit}

An apparently relevant source is M\"uller's preprint
\cite{mueller}, specifically arXiv:2104.10754v1. In its coefficient
convention, a rational $2$-function over $\Q$ is a rational series
$V(x)=\sum_{j\geq1}b_jx^j$ with integral coefficients satisfying
\begin{equation}\label{eq:two-function}
 b_{mp^r}\equiv b_{mp^{r-1}}\pmod{p^{2r}}
 \quad\text{for every prime }p\text{ and }m,r\geq1.
\end{equation}
For framing parameter $1$, Proposition 5.2(3) gives transformed
coefficients
\begin{equation}\label{eq:framing}
 f(n)=[x^n]\exp\left(n\sum_{j\geq1}\frac{b_j}{j}x^j\right).
\end{equation}
The explicit $p\geq5$ conclusion of its Theorem 1.1 asserts cubic
prime-power congruences whenever $p$ does not divide the coefficient
period. The following example satisfies these hypotheses but violates
that conclusion.

\subsection{An input that satisfies all prime conditions}

Take
\begin{equation}\label{eq:audit-input}
 V(x)=\frac{3x}{1-x}+\frac{9x^3}{1-x^3},\qquad
 b_j=3+9\,\mathbf1_{3\mid j}.
\end{equation}
The coefficient sequence has period $3$. It is integral, and
\eqref{eq:two-function} can be verified for all indices, not merely by
experiment. If $p\ne3$, multiplication by $p$ preserves divisibility
by $3$, so the two coefficients are equal. If $p=3$ and $r\geq2$,
both indices are divisible by $3$, and again the coefficients are
equal. If $p=3$ and $r=1$, the difference is either $0$ or $9$,
which is divisible by $3^2$. Thus the input is a rational $2$-function
over $\Q$ in the stated convention; no ramified-prime exception or
field extension is involved.

Formal integration gives
\[
 \sum_{j\geq1}\frac{b_j}{j}x^j
  =-3\log(1-x)-3\log(1-x^3).
\]
Hence \eqref{eq:framing} becomes
\[
 f(n)=[x^n](1-x)^{-3n}(1-x^3)^{-3n}=d_3(n).
\]
At $n=1$ and $n=5$ one obtains
\begin{align*}
 f(1)&=3,\\
 f(5)&=\binom{19}{5}+15\binom{16}{2}
       =11\,628+1\,800=13\,428,\\
 f(5)-f(1)&=13\,425=25\cdot537\equiv50\pmod{125}.
\end{align*}
The prime $5$ is unramified over $\Q$ and does not divide the period
$3$. Thus this is a counterexample to the explicit prime-by-prime
conclusion of Theorem 1.1 as printed in the cited version.

\subsection{The weighted harmonic assertion also fails}

Theorem 1.2 of the same preprint, repeated as Theorem 6.2, would require
the following period-three weighted sum at $p=n=5$ to vanish modulo
$5$ \cite{mueller}. For the very same coefficients in
\eqref{eq:audit-input}, however,
\begin{equation}\label{eq:audit-harmonic}
 \sum_{k=1}^{4}\frac{b_{5-k}b_k}{k^2}
 =9+\frac{36}{4}+\frac{36}{9}+\frac9{16}
 =\frac{361}{16}\equiv1\pmod5.
\end{equation}
The corresponding quadratic coefficient is $33/2$, which also does
not vanish modulo $5$. These are exact rational calculations.

The scope of this observation matters. A failure at one prime refutes
the cited explicit all-eligible-primes conclusion. It does not, by
itself, refute a weaker statement that allows an unspecified finite set
of exceptional primes or a compensating scalar. No such broader
refutation is asserted here.

This audit is included because otherwise it would be tempting to cite
the framing statement as a black box. The main proof does not do that.
Its logarithmic expansion and coefficient estimate are in the same
mathematical tradition, but the needed reciprocal-square statement is
proved directly for parity weights. That replacement is sufficient for
A348410 and does not inherit the false arbitrary-period assertion.

\section{Source and artifact notes}\label{app:artifacts}

Sources were inspected on 19 September 2026. The framing theorem
statements were checked against the PDF of arXiv:2104.10754v1, not only
its HTML rendering. The companion archive contains the complete source,
programs, data, and a proposed OEIS annotation for human review. No
database edit or communication with an author was performed.

\begin{thebibliography}{9}
\bibitem{oeis}
The OEIS Foundation Inc.,
\emph{A348410: Number of nonnegative integer solutions to
$n=\sum_{i=1}^{n}(a_i+b_i)$, with $b_i$ even}.
Entry initiated by C\'esar Eliud Lozada, 2021; supercongruence conjecture
contributed by Peter Bala, 21 February 2022.
\url{https://oeis.org/A348410}. Accessed 19 September 2026.

\bibitem{niu}
Tong Niu,
\emph{An explicit algebraic generating function for OEIS A348410},
arXiv:2605.16553v1, 15 May 2026.
\url{https://arxiv.org/abs/2605.16553}.

\bibitem{prodinger}
Helmut Prodinger,
\emph{The generating function of A348410 in OEIS using the diagonal
method and another sequence (A001008) from OEIS},
arXiv:2605.21255v2, 23 May 2026.
\url{https://arxiv.org/abs/2605.21255}.

\bibitem{mueller}
L. Felipe M\"uller,
\emph{Wolstenholme Type Congruences and Framing of Rational
$2$-Functions}, arXiv:2104.10754v1, 21 April 2021.
Theorem 1.1, Theorem 1.2 (also Theorem 6.2), and Proposition 5.2(3).
\url{https://arxiv.org/abs/2104.10754}.
\end{thebibliography}
\end{document}
